\documentclass[10pt]{article}

\usepackage[utf8]{inputenc}

\usepackage[T1]{fontenc}

\usepackage{amsmath}

\usepackage{amsfonts}

\usepackage{amssymb}

\usepackage[version=4]{mhchem}

\usepackage{stmaryrd}

\usepackage{graphicx}

\usepackage[export]{adjustbox}

\graphicspath{ {./images/} }



\begin{document}

ставляются иммунным множествам, хотя пммунными и продуктивными множествами не исчерпывается совокупность множеств, не являюцихся рекурсивно перечислимыми. Продуктивными оказываются многие множества, играющие важную роль в рекурсивной теории множеств (вапр., множество всех гёделевых номеров общерекурсивных функций в нек-рой гёделевой нумерации всех частично рекурсивных функций) и в ее приложениях (в тастности, множества всех номеров истинных и ложных формул әлементарной арифметики при естественной нумерации всех ее формул). Рекурсивно перечислимые множества, дополнение к-рых (до натурального ряда) является продуктивным, наз к р еа т и в н ы м и; они образуют важный класс рекурсивно перечислимых множеств.



Лит [1] Р о д ж е p c X., Теория рекурсивных функций и эффективная вычислимость, пер. с англ., М., 1972.



IPOEКTИВНАЯ АЛГЕБPА алгебра точек на проективной прямой; проективноинвариантные конструкции для определения сложения и умножения точек проективной прямой $l$, расположенной в нек-рой проективной плоскости $\pi$, для к-рой выполняется Дезарга предложение. Эти конструкции за-



\includegraphics[max width=\textwidth, center]{2023_07_08_cbd0f97ecfe4c9bba8c4g-1(1)}

висят от выбора на $l$ трех различных точек $O, E, U$.



К $о$ н с т р у кц и я I определяет для любых двух точек $A$ и $B$, отличных от $U$, третью точку $A+B$, также



Рис. 1. отличнуіо от $U$ и называемую с у мм о й т о ч е к $A$ и $B$. Для әтого в плоскости $\pi$ проводятся три прямые $a$, $b$ п $u$, отличные от $l$, не проходящие через одну точку и проходящие соответственно через точки $A, B$ и $U$. Пусть $P$ - точка пересечения прямых $u$ и $a, Q$ - точка пересечения прямых $u$ и $b, R$ - точка пересечения прямых $O Q$ и $a, S$ - точка пересечения прямых $b$ и $U R$. Тогда прямая $P S$ пересекает прямую $l$ в определенной точке $T=A+B$ (обпий случай - на рис. 1). Оказывается, так построенная точка зависит лишь от $A, B, O, U$ и не зависит от выбора прямых и точки $E$.



К о н с т р у к ц и я II определяет для любых двух точек $A$ и $B$, отличных от $U$, третью точку $A \cdot B$, также отличную от $U$, называемую п р о и 3 в е д е н и е м т о ч е к $A$ и $B$. Для әтого в плоскости проводятся три прямые $a, b, u$, проходящие соответственно через точки $A, B$ и $U$, отличные от $l$ и не проходящие через одну точку. Пусть $P$ - точка пересечения прямых $u$ и $a, Q-$ точка пересечения прямых $u$ и $b, R$ - точка пересечения прямых $E Q$ и $a, S$ - точка пересечения прямых $O R$



\begin{center}

\includegraphics[max width=\textwidth]{2023_07_08_cbd0f97ecfe4c9bba8c4g-1}

\end{center}



Рис. 2. и $b$. Тогда прямая $P S$ пересекает прямую $l$ в определенной точке $T=A \cdot B$ (общий случай на рис. 2). Оказывается, так построенная точка зависит лишь от $A, B$, $O, E, U$, но не зависит от выбора Іірямых $a, b$ и $u$.



Относит е л ь н о этих операций сложения и умноже-



ния точки прямой $l$ (отличные от $U)$ образуют тело $K(O$, $E, U)$. Поменяв ролями $A$ и $B$ в конструкции II, получают инверсно изоморфное тело $K^{*}(O, E, U)$. Если $O^{\prime}$, $E^{\prime}, U^{\prime}$ - любая другая упорядоченная тройка точек на прямой $l$ из той же плоскости $\pi$, то соответствующее тело $K^{\prime}\left(O^{\prime}, E^{\prime}, U^{\prime}\right)$ изоморфно $K(O, E, U)$ вследствие того, что между прямыми $l$ и $l^{\prime}$ существует проективное соответствие; поэтому любое тело $K$, изоморфное им, наз. просто телом данной проективной плоскости (или даже данной проективной геометрии); говорят также, что имеет место п р о е к т и в на я ге о ме т р и я н а д т е л о м $K$. В общих случаях конструкций I и II фигурируют четыре лежащие в одной плоскости точки $P, Q$, $R, S$, никакие три из к-рых не коллинеарны; они образуют т. н. (полный) ч е ты р е х в е р ши и ни к с тремя парами противоположных сторон $P Q, R S ; P S$, $Q R ; P R, Q S$. Точки пересечения $Z, X, Y$ этих пар противоположных сторон называется д и а го н а л ьн ы м и т о ч к а м и, а прямые, соединяющие диагональные точки, - д и а г о н а л я м и. Специальный случай, не показанный на рисунке, соответствует той ситуации, когда $X, Y, Z$ коллинеарны (см. Фано поcmулат).



Аналогичные построения проводятся и в пучке прямых, проходящих через одну точку с использованием (полного) ч е т ы р е х с т о р о н н и к а - фигуры, двойственной четырехвершиннику, и приводят к телу $K^{*}$, инверсно изоморфному $K$.



Свойства проективной прямой $l$ как алгебраич. системы определяются геометрическими (проективно-инвариантными) свойствами проективной плоскости, в к-рой она расположена. Так, напр., коммутативность $K$ әквивалентна выполнению Паппа аксиомы; то, что характеристика $K$ не равна 2 , эквивалентно постулату Фано; при отсутствии автоморфизмов у тела $K$, отличных от внутренних, любое проективное преобразование есть коллинеачия, и т. Д.



С помощыо тела $K$ на прямой, а затем и в проективном пространстве, ее содержащем, вводятся проективные координаты, описывающие алгебраич. модель проективного пространства, так что содержание проективной геометрии по существу определяется свойствами того тела $K$, над к-рым она построена.



В широком смысле П. а. исследует совокупность подпространств проективного пространства, являющуюся дедекиндовой решеткой с дополнениями; при әтом конечномерности пространства не требуется, но накладываются условия полноты, существования однородного базиса и т. д., благодаря чему устанавливаются разнообразные связи с теорией простых и регулярных колец, теорий операторных абелевых групп и др. разделами алгебры.



Лum.: [1] Х о д Ж В., П и д о Д., Методы алгебраической геометрии, пер. с англ., т. 1, М., 1954; [2] А р т и н Э., Гео-

метрическая алгебра, пер. с англ., М., 1969.



М. И. Войчеховский. IРОЕКТИВНАЯ ГЕОМЕТРИЯ - раздел геометрии, изучающий свойства фигур, не меняющиеся при проективных преобразованиях, напр при проектировании. Такие свойства наз. п р о е к т и в ны м и; к ним относятся, напр., прямолинейное расположение точек (к о л ли н е а р н о с т ь), порядок алгебраич. кривой и т. Д.



При проектировании точек одной плоскости П на другую плоскость $\Pi^{\prime}$ не каждая точка П' имеет прообраз в П и не каждая точка из П имеет образ в $\Pi^{\prime}$. Это обстоятельство привело к необходимости дополнения евклидова пространства т. н. бесконечно удаленныли элементами (н е с об с т в е н ны м и точк а м и, п р я м м и и п л о с к о с т в ю) и кобразованию нового геометрич. объекта - трехмерного проективного пространства. При этом каждая прямая дополняется одной несобственной точкой, каждая плоскость - одной несобственной прямой, все пространство - одной несобственной плоскостью. Параллельные прямые дополняются одной и той же несобственной точкой, непараллельные - разными, параллельные плоскости





\end{document}